\section{Finite-sample cost}
\label{sec:finite_sample}

Suppose the population coefficient aligns with $(1,-1)$, so $D$ is linearly sufficient at the population level.
A correctly specified probit in $(A,B)$ still fits one more coefficient than a probit in $D$.

\begin{proposition}[Parameter-saving gain]
\label{prop:gain}
For any correctly specified regular model that differs by one free parameter,
\[
  \mathbb{E}\bigl[\mathrm{LL}_{\mathrm{full}}-\mathrm{LL}_D\bigr]
  = \frac{1}{2n}+O(n^{-2}),
\]
where $\mathrm{LL}$ denotes expected test log-loss and $n$ is the training sample size~\parencite{akaike1973,stone1977}.
Under joint normality with $\beta\propto(1,-1)$, a probit with intercept is correctly specified for both feature sets, so the expansion applies.
\end{proposition}

The leading term is the usual one-parameter Akaike penalty.
On three aligned Gaussian designs A1--A3, a weighted fit of $n\cdot\mathrm{gain}=a+\gamma/n$ gave leading coefficients $0.491$, $0.503$ and $0.489$ (standard errors $0.007$).
The second-order coefficients were $17.3$, $16.6$ and $15.6$.
The pooled fitted value is $\gamma\approx 16.6$.
That coefficient is an empirical fit on those three designs.
It is not claimed to be universal.

When $D$ is slightly misaligned, the population gap $\Delta_\infty=\mathrm{LL}_{D,\infty}-\mathrm{LL}_{\mathrm{full},\infty}$ is positive and computable from the Gaussian $(z,r^2)$ map.
The finite-sample advantage of $D$ is then approximately $1/(2n)-\Delta_\infty$.
The first-order crossover $n^*\approx 1/(2\Delta_\infty)$ is the sample size at which the expected likelihood-ratio statistic for the redundant coefficient equals Akaike's threshold of two.

\begin{remark}[Crossover approximation]
\label{rem:crossover}
If $\Delta_\infty>0$, the first-order crossover sample size is $n^*\approx 1/(2\Delta_\infty)$.
With the second-order remainder and the fitted coefficient $\gamma$ from the aligned designs,
\[
  n^*
  = \frac{\tfrac12+\sqrt{\tfrac14+4\gamma\Delta_\infty}}{2\Delta_\infty}.
\]
The formula reuses the aligned $\gamma$ for misaligned cells.
That transfer is an assumption, not a derived identity.
\end{remark}

\Cref{tab:crossover} compares the two predictions with the empirical crossover on four cells.
Cell~B misaligns through the outcome coefficient ratio $b_B/b_A$.
Cell~C misaligns through noise asymmetry $\kappa_N\neq 1$ at $b=c$.
The second-order formula uses the pooled $\gamma$ from A1--A3, so the B and C rows are out-of-sample for that coefficient.
For B with nominal $n^*=1000$, the empirical crossover $980.1$ (interval $950$--$1023$) is closer to the first-order value $1000$ than to the second-order value $1032.2$; the gap is within Monte Carlo noise.
The expansion uses the Gaussian $(z,r^2)$ map.
It was not re-estimated on the stress laws of~\cref{sec:stress}.

\begin{table}[t]
  \centering
  \caption{Crossover sample sizes.
           Second-order predictions use the pooled fitted $\gamma\approx 16.6$ from the aligned designs A1--A3.
           Empirical intervals are from the E3 summary.}
  \label{tab:crossover}
  \begin{tabular}{@{}lccccc@{}}
    \toprule
    Cell & $\Delta_\infty$ & First-order & Second-order & Empirical & Interval \\
    \midrule
    B, $b$-ratio, $n^*=100$  & $0.005$    & $100$  & $126.3$  & $124.8$ & $120.8$--$128$ \\
    B, $b$-ratio, $n^*=300$  & $0.001667$ & $300$  & $330.2$  & $327.6$ & $309.9$--$340.9$ \\
    B, $b$-ratio, $n^*=1000$ & $0.0005$   & $1000$ & $1032.2$ & $980.1$ & $950$--$1023$ \\
    C, $\kappa_N$, $n^*=300$ & $0.001667$ & $300$ & $330.2$ & $342$ & $327.4$--$353.8$ \\
    \bottomrule
  \end{tabular}
\end{table}
